melodi - Shu-bi-dua's rendition af Midsommersangen (Vi elsker vort land)

\[Vers 1\] # Skal handle om at det er 50 års jubilæum
Vi elsker vort klub, nu fem-ti år som det går
Har vi samlet os her med glæde for øje
Når vi mindes de år, F-embererne de har budt
Lad os hylde og hilse de ny' der fornøje
Vi synger vort lov over vej, over gade
Vi kranser vort navn, her til fest i en "lade"

\[Omkvæd 1, version 1\]
|: For den skønneste krans, bli'r da' din, F-klubben!
Den er bunden af F-embers hjerter så varme, så glade :|

\[Omkvæd 1, version 2\]
|: I det skønneste sted, hør' du til F-klubben!
I bunden af F-embernes hjerte så varme, så glade :|

\[Omkvæd 1, version 3\] høhø
|: I det skønneste sted, hør' du til F-klubben!
I bunden af F-embernes maver så varme, så glade :|

\[Vers 2\] # Skal handle om hvad foreningen bidrager/betyder nu
== Vi elsker vort klub, mens på uni vi går/gik
Den ligger som om det var skabt til at ......==

\[OLD Vers 2\]
Vi elsker vort land, men ved midsommer mest
Når hver sky over marken velsignelsen sender
Når af blomster er flest, og når kvæget i spand
Giver rigeligst gave til flittige hænder
Når ikke vi pløjer og harver og tromler
Når koen sin middag i kløveren gumler



\[Omkvæd 2\]
|: Da går F-embers til dans  
På dit bud F-Klubben!  
==Ret som føllet og lammet, der frit over engen sig tumler== :|


\[Vers 3\] # Skal handle om Foreningens fremtid / overdragelse ad pinden

\[OLD Vers 3\]
Vi elsker vort land, og med sværdet i hånd  
Skal være udenvælts fjende beredte os kende
Men mod ufredens ånd, over mark, over strand
Vil vi bålet på fædrenes gravhøje tænde
Hver by har sin heks, og hver sogn sine trolde
Dem vil vi fra livet med glædesblus holde

\[Omkvæd 3\]
|: Der er fest, der er dans F-Klubben, F-Klubben!  
==Den kan vindes, hvor hjerterne aldrig bli'r tvivlende kolde! :|==